% Generated by render_latex.py. Draft with explicit evidence placeholders.
\documentclass[sigconf,anonymous]{acmart}
\usepackage{amsmath,booktabs,tabularx,array,listings,tikz}
\usetikzlibrary{arrows.meta,positioning}
\setlength{\emergencystretch}{1em}
\setcopyright{none}
\settopmatter{printacmref=false}
\acmDOI{}
\acmISBN{}
\acmConference[SIGMOD submission draft]{}{}{}
\lstset{basicstyle=\ttfamily\small,breaklines=true,columns=fullflexible,keepspaces=true,frame=single}
\title{ExRaBitQ-Disk: Controlling Payload I/O in Quantized Graph Search}
\begin{document}
\begin{abstract}
Graph-based approximate nearest neighbor search must balance retrieval accuracy, memory consumption, and storage access when an index exceeds its DRAM budget. Compact vector codes reduce representation size, but a candidate requiring only a small distance computation can still trigger a full disk-page read. We present ExRaBitQ-Disk, a quantized graph search system that uses a resident one-bit representation to decide whether a candidate warrants access to its disk-resident payload. The system constructs a Vamana graph using symmetric four-bit distance estimates and performs online search in three stages: one-bit candidate screening, four-bit traversal scoring, and residual-corrected reranking. Screening and traversal share a query representation and a partial inner product; page coalescing and bounded query-local reuse translate surviving candidates into physical reads. We analyze the approximation errors of construction and querying, and distinguish candidate filtering from page savings. Our evaluation protocol separates fixed-candidate quantizer measurements, controlled graph and storage ablations, and complete-system comparisons under a common memory budget. \textbf{[E1: Insert audited, recall-matched throughput and physical-I/O results across the admitted datasets.]} The design treats quantized representations as a means of controlling payload access as well as reducing arithmetic and storage costs.
\end{abstract}
\keywords{Approximate nearest neighbor search, vector quantization, graph indexes, disk I/O}
\maketitle
\newcommand{\diskfigure}[1]{\input{figures/disk_w32/#1.tex}}
\input{body}
\bibliographystyle{ACM-Reference-Format}
\bibliography{references}
\end{document}
